\chapter{Bức Họa Nhân Sinh Toàn Vẹn (The Complete Portrait of Life)}
\begin{center}
  \textit{\color{truyengold}``Một cuộc đời trọn vẹn không phải là cuộc đời không có đau khổ --- mà là cuộc đời đủ sâu để chứa đựng cả đau khổ lẫn hạnh phúc.''}\\[4pt]
  \textit{(A complete life is not a life without suffering --- but a life deep enough to hold both suffering and joy.)}\\[4pt]
  \small\color{truyenblue}--- Rainer Maria Rilke ---
\end{center}
\ngancach

\section{Thọ 100 Tuổi --- Nhìn Lại Cuộc Đời}
\begin{truyen}{Thọ 100 Tuổi --- Nhìn Lại Cuộc Đời}{Câu chuyện Triết lý}
\chuhoa{M}{ột} cụ già 100 tuổi được phóng viên hỏi: ``Nếu được sống lại, cụ muốn thay đổi điều gì?''\\[4pt]
\textit{(A 100-year-old elder was asked by a journalist: `If you could live again, what would you change?')}

Cụ suy nghĩ lâu rồi trả lời: ``Tôi sẽ lo lắng ít hơn --- vì 90\% những thứ tôi lo lắng không bao giờ xảy ra. Tôi sẽ yêu thương nhiều hơn và sẽ cởi mở hơn với những sai lầm của người khác. Và tôi sẽ đi xem thêm nhiều hoàng hôn hơn.''\\[4pt]
\textit{(The elder thought long and then replied: `I would worry less --- because 90\% of what I worried about never happened. I would love more and be more open to others' mistakes. And I would watch more sunsets.')}

Phóng viên hỏi: ``Còn điều gì cụ hài lòng nhất?'' Cụ mỉm cười: ``Rằng tôi không bao giờ bỏ cuộc --- dù vấp ngã bao lần, tôi vẫn đứng dậy.''\\[4pt]
\textit{(The journalist asked: `What are you most satisfied with?' The elder smiled: `That I never gave up --- no matter how many times I fell, I always got back up.')}

\end{truyen}
\begin{baihoc}
  Cuộc đời không cần phải hoàn hảo --- chỉ cần được sống hết mình, yêu hết lòng, và không bao giờ bỏ cuộc.\\[4pt]
  \textit{(Life does not need to be perfect --- it only needs to be lived fully, loved wholeheartedly, and never given up on.)}
\end{baihoc}

\section{Bức Thư Gửi Con Của Albert Einstein}
\begin{truyen}{Bức Thư Gửi Con Của Albert Einstein}{Albert Einstein --- Hoa Kỳ, 1915}
\chuhoa{N}{ăm} 1915, Albert Einstein viết thư cho con trai Hans Albert đang học tại Zurich, gửi vào trong đó những lời dặn dò mà ông nghĩ là quan trọng nhất.\\[4pt]
\textit{(In 1915, Albert Einstein wrote to his son Hans Albert studying in Zurich, including in that letter what he considered the most important advice.)}

Ông không viết về vật lý, không viết về công thức hay lý thuyết --- ông viết về tình yêu.\\[4pt]
\textit{(He did not write about physics, formulas, or theories --- he wrote about love.)}

``Con hãy biết rằng đó là cách ba học tốt nhất, và khi ba đã già --- ba vẫn học được nhiều nhất từ những gì ba yêu thích nhất.''\\[4pt]
\textit{(`Know that this is how Father learns best, and even now that Father is older --- he still learns most from what he loves most.')}

Người tìm ra lý thuyết tương đối của vũ trụ tin rằng bí quyết của mọi khám phá lớn chỉ là một từ: tình yêu --- yêu điều mình làm.\\[4pt]
\textit{(The man who discovered the theory of relativity of the universe believed the secret of every great discovery is one word: love --- loving what you do.)}

\end{truyen}
\begin{baihoc}
  Thiên tài không đến từ trí tuệ vượt trội --- mà từ tình yêu đủ sâu để kiên trì theo đuổi điều mình tin.\\[4pt]
  \textit{(Genius does not come from superior intellect --- but from love deep enough to persistently pursue what one believes in.)}
\end{baihoc}

\section{Lời Trăng Trối Của Tolstoy}
\begin{truyen}{Lời Trăng Trối Của Tolstoy}{Leo Tolstoy --- Nga, 1910}
\chuhoa{T}{rước} khi qua đời, nhà văn vĩ đại Leo Tolstoy --- tác giả của \textit{Chiến Tranh Và Hòa Bình} và \textit{Anna Karenina} --- thốt lên những lời cuối: ``Tôi không hiểu gì cả.''\\[4pt]
\textit{(Before dying, the great writer Leo Tolstoy --- author of \textit{War and Peace} and \textit{Anna Karenina} --- uttered his last words: `I understand nothing.')}

Người viết nên những cuốn sách đồ sộ nhất về con người, xã hội và đạo đức --- đến cuối đời vẫn thừa nhận mình chưa hiểu hết cuộc sống.\\[4pt]
\textit{(The man who wrote the most monumental books about people, society, and ethics --- at the end of his life still admitted he had not fully understood life.)}

Đây không phải là thất bại --- đây là sự khôn ngoan thật sự: người càng học nhiều, càng hiểu mình còn biết ít đến nhường nào.\\[4pt]
\textit{(This was not failure --- this was true wisdom: the more one learns, the more one understands how little one knows.)}

Socrates từng nói: ``Điều duy nhất tôi biết là tôi không biết gì cả'' --- và đó chính là khởi đầu của mọi trí tuệ.\\[4pt]
\textit{(Socrates once said: `The only thing I know is that I know nothing' --- and that is the very beginning of all wisdom.)}

\end{truyen}
\begin{baihoc}
  Khiêm tốn trước sự bí ẩn vô tận của cuộc sống là nền tảng của mọi hiểu biết thật sự.\\[4pt]
  \textit{(Humility before the infinite mystery of life is the foundation of all true understanding.)}
\end{baihoc}

\section{Lời Chào Cuối --- Hành Trình Tiếp Tục}
\begin{truyen}{Lời Chào Cuối --- Hành Trình Tiếp Tục}{Tuyển Tập Điển Tích \& Chuyện Kể --- Quyển I đến VI}
\chuhoa{B}{ạn} vừa đọc qua sáu quyển sách --- hàng trăm câu chuyện, hàng trăm bài học --- từ cổ đại đến hiện đại, từ Đông sang Tây, từ thiên nhiên đến con người.\\[4pt]
\textit{(You have just read through six volumes --- hundreds of stories, hundreds of lessons --- from ancient to modern, from East to West, from nature to humanity.)}

Mỗi câu chuyện là một hạt giống. Hạt giống không nảy mầm trong lúc đọc --- mà nảy mầm trong cuộc sống thực, khi bạn gặp thử thách, khi bạn phải chọn lựa, khi bạn cần sức mạnh.\\[4pt]
\textit{(Each story is a seed. Seeds do not sprout while reading --- they sprout in real life, when you face challenges, when you must choose, when you need strength.)}

Hành trình không kết thúc ở trang cuối cuốn sách này --- nó chỉ bắt đầu trong trái tim của người đọc.\\[4pt]
\textit{(The journey does not end on the last page of this book --- it only begins in the heart of the reader.)}

Câu hỏi cuối cùng của bộ sách này không phải là ``Bạn đã học được gì?'' --- mà là ``Bạn sẽ sống như thế nào?''\\[4pt]
\textit{(The final question of this series is not `What have you learned?' --- but `How will you live?')}

\end{truyen}
\begin{baihoc}
  Đọc để hiểu --- Hiểu để sống --- Sống để truyền lại cho thế hệ tiếp theo.\\[4pt]
  Đó là vòng tròn của tri thức và tình người không bao giờ ngừng.\\[8pt]
  \textit{(Read to understand --- Understand to live --- Live to pass on to the next generation.)}\\
  \textit{(That is the eternal circle of knowledge and humanity.)}
\end{baihoc}
